% GENERATED FILE. Edit canonical lessons, then run npm run generate:books.
% canonical-source-hash: fnv1a64:f6b74168ca7d8e0d
% canonical-lessons: ML-C318-parajayam, ML-C318-telivu, ML-C318-vicaram, ML-C318-ormma, ML-C318-nirddesam

\chapter{Failure, Evidence, Thought, Memory, Suggestion}
\label{ch:ml-failure-evidence-thought-memory-suggestion}
\hlchaptermodality{318}

\begin{chapteropening}
\textbf{By the end of this chapter:} \emph{I can say failure, evidence, a thought, a memory and a suggestion.}

Complete the last lesson of chapter 318: I can say failure, evidence, a thought, a memory and a suggestion.
\end{chapteropening}

\section[\texorpdfstring{\ml{പരാജയം} (parājayaṁ)}{പരാജയം (parājayaṁ)}]{\ml{പരാജയം} (parājayaṁ) — failure}

\label{lesson:ML-C318-parajayam}



\begin{quote}
Before the new one: say the Malayalam for not at all, then the Malayalam for barely.
\end{quote}

\subsection*{You'll want to know: \ml{പരാജയം}}
\textbf{\ml{പരാജയം}} — \emph{parājayaṁ} — \textquotedblleft{}failure\textquotedblright{}.

\begin{grammarlens}[title={What you've built}]
One more word for this chapter.
\end{grammarlens}

\subsection*{Your turn}
\begin{itemize}
\raggedright
  \item \emph{Say it:} \emph{parājayaṁ}
  \item \emph{Say it:} \emph{parājayaṁ}, once more
  \item \emph{Say it:} say \emph{parājayaṁ}
  \item \emph{Recall:} say the Malayalam for not at all, then the Malayalam for barely, then say \emph{parājayaṁ} again
\end{itemize}

\begin{tcolorbox}[breakable,colback=teal!4,colframe=teal!35!black,title={Before you move on}]
What does \ml{പരാജയം} mean? (\textquotedblleft{}Failure\textquotedblright{}.) Say it once more.
\end{tcolorbox}
\section[\texorpdfstring{\ml{തെളിവ്} (teḷivŭ)}{തെളിവ് (teḷivŭ)}]{\ml{തെളിവ്} (teḷivŭ) — evidence, proof}

\label{lesson:ML-C318-telivu}



\begin{quote}
Before the new one: say the Malayalam for barely, then the Malayalam for failure.
\end{quote}

\subsection*{You'll want to know: \ml{തെളിവ്}}
\textbf{\ml{തെളിവ്}} — \emph{teḷivŭ} — \textquotedblleft{}evidence, proof\textquotedblright{}.

\begin{grammarlens}[title={What you've built}]
One more word for this chapter.
\end{grammarlens}

\subsection*{Your turn}
\begin{itemize}
\raggedright
  \item \emph{Say it:} \emph{teḷivŭ}
  \item \emph{Say it:} \emph{teḷivŭ}, once more
  \item \emph{Say it:} say \emph{teḷivŭ}
  \item \emph{Recall:} say the Malayalam for barely, then the Malayalam for failure, then say \emph{teḷivŭ} again
\end{itemize}

\begin{tcolorbox}[breakable,colback=teal!4,colframe=teal!35!black,title={Before you move on}]
What does \ml{തെളിവ്} mean? (\textquotedblleft{}Evidence, proof\textquotedblright{}.) Say it once more.
\end{tcolorbox}
\section[\texorpdfstring{\ml{വിചാരം} (vicāraṁ)}{വിചാരം (vicāraṁ)}]{\ml{വിചാരം} (vicāraṁ) — a thought}

\label{lesson:ML-C318-vicaram}



\begin{quote}
Before the new one: say the Malayalam for failure, then the Malayalam for evidence.
\end{quote}

\subsection*{You'll want to know: \ml{വിചാരം}}
\textbf{\ml{വിചാരം}} — \emph{vicāraṁ} — \textquotedblleft{}a thought\textquotedblright{}.

\begin{grammarlens}[title={What you've built}]
One more word for this chapter.
\end{grammarlens}

\subsection*{Your turn}
\begin{itemize}
\raggedright
  \item \emph{Say it:} \emph{vicāraṁ}
  \item \emph{Say it:} \emph{vicāraṁ}, once more
  \item \emph{Say it:} say \emph{vicāraṁ}
  \item \emph{Recall:} say the Malayalam for failure, then the Malayalam for evidence, then say \emph{vicāraṁ} again
\end{itemize}

\begin{tcolorbox}[breakable,colback=teal!4,colframe=teal!35!black,title={Before you move on}]
What does \ml{വിചാരം} mean? (\textquotedblleft{}A thought\textquotedblright{}.) Say it once more.
\end{tcolorbox}
\section[\texorpdfstring{\ml{ഓർമ്മ} (ōrmma)}{ഓർമ്മ (ōrmma)}]{\ml{ഓർമ്മ} (ōrmma) — a memory}

\label{lesson:ML-C318-ormma}



\begin{quote}
Before the new one: say the Malayalam for evidence, then the Malayalam for a thought.
\end{quote}

\subsection*{You'll want to know: \ml{ഓർമ്മ}}
\textbf{\ml{ഓർമ്മ}} — \emph{ōrmma} — \textquotedblleft{}a memory\textquotedblright{}.

\begin{grammarlens}[title={What you've built}]
One more word for this chapter.
\end{grammarlens}

\subsection*{Your turn}
\begin{itemize}
\raggedright
  \item \emph{Say it:} \emph{ōrmma}
  \item \emph{Say it:} \emph{ōrmma}, once more
  \item \emph{Say it:} say \emph{ōrmma}
  \item \emph{Recall:} say the Malayalam for evidence, then the Malayalam for a thought, then say \emph{ōrmma} again
\end{itemize}

\begin{tcolorbox}[breakable,colback=teal!4,colframe=teal!35!black,title={Before you move on}]
What does \ml{ഓർമ്മ} mean? (\textquotedblleft{}A memory\textquotedblright{}.) Say it once more.
\end{tcolorbox}
\section[\texorpdfstring{\ml{നിർദ്ദേശം} (nirddēśaṁ)}{നിർദ്ദേശം (nirddēśaṁ)}]{\ml{നിർദ്ദേശം} (nirddēśaṁ) — a suggestion, an instruction}

\label{lesson:ML-C318-nirddesam}



\begin{quote}
Before the new one: say the Malayalam for a thought, then the Malayalam for a memory.
\end{quote}

\subsection*{You'll want to know: \ml{നിർദ്ദേശം}}
\textbf{\ml{നിർദ്ദേശം}} — \emph{nirddēśaṁ} — \textquotedblleft{}a suggestion, an instruction\textquotedblright{}.

\begin{grammarlens}[title={What you've built}]
One more word for this chapter.
\end{grammarlens}

\subsection*{Your turn}
\begin{itemize}
\raggedright
  \item \emph{Say it:} \emph{nirddēśaṁ}
  \item \emph{Say it:} \emph{nirddēśaṁ}, once more
  \item \emph{Say it:} say \emph{nirddēśaṁ}
  \item \emph{Recall:} say the Malayalam for a thought, then the Malayalam for a memory, then say \emph{nirddēśaṁ} again
\end{itemize}

\begin{tcolorbox}[breakable,colback=teal!4,colframe=teal!35!black,title={Before you move on}]
What does \ml{നിർദ്ദേശം} mean? (\textquotedblleft{}A suggestion, an instruction\textquotedblright{}.) Say it once more.
\end{tcolorbox}
